\begin{lemma}\label{lem:decic-marked-two-cycle}
Let $K$ be a number field of degree $10$, let $E/K$ satisfy
$\operatorname{End}_K(E)=\mathbb Z$, and let $P\in E(K)$ have
odd order $n$. Suppose that neither $2$ nor $32$ is congruent to
$1$ or $-1$ modulo $n$. There is no degree-$2$ isogeny
$\psi:E\to E^\sigma$, with $\sigma\in\operatorname{Aut}(K)$,
such that $\psi(P)=aP^\sigma$ for $a=1$ or $2$.
\end{lemma}
\begin{proof}
Let $h$ be the order of $\sigma$, so $h\mid10$.
Composing the $h$ conjugate isogenies gives a $K$-endomorphism
of degree $2^h$ which acts on $P$ by $a^h$.
It is multiplication by an integer. Hence $h$ is even, so
$h=2$ or $10$, and the multiplier is $\pm2^{h/2}$.
For either value of $a$, comparison on $P$ gives
$2^{h/2}\equiv\pm1\pmod n$, contrary to the hypothesis.
\end{proof}

\begin{proposition}\label{prop:decic-no90}
An elliptic curve over a number field of degree $10$ has no point
of order $90$.
\end{proposition}
\begin{proof}
Suppose that $E/K$ has a point $Q$ of order $90$, with
$[K:\mathbb Q]=10$, and put $P=2Q$.
Let $D$ be the degree-$10$ rational trace divisor of $(E,P)$
on $X=X_1(45)$, of genus $41$.
The rational Jacobian has rank zero by
\cite[Theorem~3.1]{DEHMZ}, and specialization at $2$ is injective
on it. We recall the degree-independent injectivity argument.
Its rational $2$-primary torsion is $C_4\times C_{32}$, whose
automorphism group is a $2$-group.
Thus the order-$3$ diamond subgroup $\{1,16,31\}$ acts trivially.
Over $F=\mathbb Q(\zeta_3)$ its quotient is $X_1(3,15)$.
The sixteen rational cusps over $\mathbb F_4$ lift to $F$,
and their $2$-primary specialization image has order $8192$.
The independent good reduction over $\mathbb F_7$ gives the same
upper bound for all $F$-rational $2$-primary torsion.
Hence target specialization is injective; the norm identity
$1+h+h^2=[3]$ on the source rational $2$-primary group proves
source injectivity. Odd torsion is automatically injective.
The exact cusp labels and both finite class-group computations
are retained with their successful native exit records.
We also use
\[
                  \operatorname{gon}_{\mathbb Q}X=18,
\]
by \cite[Theorem~D]{Varivoda26}.
At each odd good prime $q$, all rational torsion, including its
$q$-primary part, injects under reduction by the unramified
formal-group lemma. Thus Eichler--Shimura gives
$(T_q-1-q\langle q\rangle)J_X(\mathbb Q)=0$.

The complete noncuspidal place counts at $2$ through degree $10$
are
\[
                       (0,0,0,0,0,0,0,60,60,108).
\]
A reduction containing a degree-$8$ or degree-$9$ place has a
cuspidal complement of degree $2$ or $1$, respectively.
Every such complement lifts globally.
The accepted exact Hecke nonprincipality tests exclude all five
orbits in either degree: in degree $9$ all five tests use
$T_7-1-7\langle7\rangle$.
They apply with the present complement because the rational
Jacobian is finite and the complete rational cusp blocks are
annihilated by the cusp boundary formula.

We next exclude entirely cuspidal reduction.
The finite cusp signature is $1^{12}2^4 4^5 6^2 12$.
Globally the degrees are
\[
                     1^{12}2^4 4^3 6^2\,8\,12.
\]
Thus every reduction containing no quartic cusp lifts to a
rational cuspidal divisor $C$.
Injectivity gives $D\sim C$, and gonality gives $D=C$,
a contradiction.
For reductions containing a quartic cusp, use
$U=(\mathcal C-1)H^0(X_{\mathbb F_2},\Omega^1)$, of dimension $39$.
All required restrictions occur in the accepted degree-$8$
local data: after a quartic cusp is selected, the remaining
degree is at most $6$.
Independent incremental and complete-matrix computations check
all $94420$ degree-$10$ configurations containing a quartic cusp.
Every restriction has rank $10$.
The full-Jacobian dyadic criterion, using the entire closed-cusp
trace over $\mathbb Q_2$, therefore excludes these reductions
as well. No quotient cotangent exactness at $2$ is used.

Consequently $\overline D$ is a single noncuspidal degree-$10$
place. The moduli field is $K$, and there is a unique prime over
$2$, unramified with residue field $\mathbb F_{1024}$.
The fine level-$45$ model gives good reduction of $E$ there.
The complete independent elliptic-curve census gives orders
$990$ or $1080$. Both reductions are ordinary.

We claim that $\operatorname{End}_K(E)=\mathbb Z$.
Otherwise the CM field $L$ embeds in $K$, by its action on the
invariant differential. Ordinary specialization identifies $L$
with the quadratic field generated by Frobenius.
The two possible traces, $35$ and $-55$, give
\[
                         L=\mathbb Q(\sqrt{-319})
                  \quad\text{or}\quad\mathbb Q(\sqrt{-119}).
\]
Both fields have class number $10$, as checked by complete reduced
binary quadratic forms and independently by PARI.
Every order in either field has class number at least $10$,
by the ring class number formula \cite[Section~7]{Cox}.
The main theorem of complex multiplication
\cite[Section~11]{Cox} therefore gives
$[L(j(E)):L]\geq10$.
However, $L(j(E))\subseteq K$ and $[K:L]=5$, a contradiction.

Let $R=45Q$ and let $\phi:E\to E'=E/\langle R\rangle$ be the
$K$-defined $2$-isogeny. Put $P'=\phi(P)$, of order $45$, and let
$D'$ be its degree-$10$ trace on $X$.
The isogeny extends over the good elliptic models at $2$.
On the ordinary special fibre its kernel is either connected
or etale. Thus its marked image is Frobenius, or Verschiebung
with the level multiplied by $2$.
Since the closed divisor $\overline D$ is Frobenius invariant,
\[
                   \overline {D'}=\overline D
          \quad\text{or}\quad
                   \overline {D'}=\langle2\rangle\overline D.
\]
The residue Frobenius is an automorphism, so no extra norm
multiplicity occurs.
Injectivity and gonality $18>10$ give
$D'=D$ or $D'=\langle2\rangle D$.
Both divisors are closed of degree $10$.
The fine moduli interpretation therefore identifies the marked
isogeny, after a $K$-isomorphism, with
\[
             \psi:E\longrightarrow E^\sigma,
             \qquad\psi(P)=P^\sigma\quad\text{or}\quad2P^\sigma,
\]
for some $\sigma\in\operatorname{Aut}(K)$.
Lemma~\ref{lem:decic-marked-two-cycle} applies with $n=45$,
since $2,32\not\equiv\pm1\pmod{45}$, and gives a contradiction.
\end{proof}
